\section{Lean formalization}\label{app:formalization}
\subsection*{Trust base and methodology}

The Lean development for this paper is deliberately mathlib-free.  It
does not import an external mathematics library, and therefore no
mathlib axioms occur in the dependency closure of the BSD declarations.
Every declaration entered in the manifest with theorem status was
audited with \texttt{\#print axioms}.  The accepted trust base is
Lean's standard foundational base: \texttt{propext},
\texttt{Quot.sound}, and \texttt{Classical.choice}, with
\texttt{funext} and the \texttt{choice} recursor appearing as standard
consequences in the audit output.  No project-local axiom lies behind
any theorem-marked result.  The source tree also contains no active
\texttt{axiom}, \texttt{opaque}, \texttt{constant}, \texttt{sorry}, or
\texttt{admit} declarations in the BSD namespace.  This point matters
for the interpretation of the recognition sources and imports: they
are represented as typed structure carriers, never as Lean axioms.

The formalization mechanizes the framework-structural closure
architecture, not the imported arithmetic geometry.  There are three
main disclosure modes.  First, typed structure carriers are used for
the imported theorem stacks and the recognition sources.  The
Cassels--Tate, Selmer-complex pairing, Beilinson, overconvergent
\(\eta\), and cascade imports are records whose mathematical content is
carried as hypothesis fields.  The three recognition sources are
likewise records: their arithmetic content is carried explicitly and
then consumed by later statements.  In all of these cases the content
is carried, not derived, and it is not an axiom of the formalization.

Second, several results are projection-packaged conditional schemas.
The \(\chiCTp\) comparison, the overconvergent \(\eta\)-formula, the
rank-\(\leq 1\) cascade theorem, the composite signature, and the
headline conditional landing are checked projections over explicit
proof-carrying records, conditional on the relevant imports and
recognition sources.  The formalization checks that the projections and
implication chains use the supplied records in the stated way; it does
not independently derive the underlying arithmetic theorems.  For this
reason the body footnotes say that these claims are tracked as
conditional schemas, not that the formalization derives the arithmetic
comparison unconditionally.

Third, some statements are substantive finite Lean proofs of the
structural layer.  The readout-level equivalence, the
master-theorem-applicability statement, and the AOR mechanical and
recognition discharge theorems are proofs about the finite record
architecture and its bookkeeping.  The AOR-instance membership is
checked over a narrowed representation: the local syntactic
refinement-stable predicate used by this paper, not the full Audited
Operational Realisability meta-theory.  Thus Lean checks the closure
architecture---the residual readout algebra, the implication chains,
and the discharge bookkeeping---while certifying none of the imported
arithmetic geometry or the recognition-source content itself.

In coverage terms, the paper has 28 statements of record.  Of these,
21 are mechanized: 4 substantive structural theorems, 5
projection-packaged conditional schemas, 8 typed definition carriers
(including the shell and predicate mirrors, the typed import carriers,
and the AOR-instance object), 3 recognition-source carriers, and 1
narrowed AOR-instance membership.  The remaining 7 are not mechanized:
the two audit-result no-gos, two remarks, and three non-claims.  The
tables below give the per-label coverage and the corresponding Lean
declaration map.
\subsection*{Wording-discipline table}

Every mechanized statement in the body carries a footnote whose wording
is canonical, not editorial.  The wording encodes whether the result is
a substantive finite Lean proof, a projection-packaged conditional
schema, a narrowed-representation check, a typed structure carrier, or
a recognition-source carrier.  Projection-packaging and narrowing are
therefore visible at the point of disclosure; unmechanized statements
carry no footnote.

\begin{center}
\begin{tabularx}{\textwidth}{@{}>{\raggedright\arraybackslash}p{0.36\textwidth}X@{}}
\toprule
Statement class & Footnote convention \\
\midrule
Theorem, substantive finite proof &
Verified in Lean as \texttt{[decl]}; see App~\ref{app:formalization}. \\
\addlinespace
Theorem, projection-packaged conditional schema &
Tracked in the formalization harness as a conditional schema
\texttt{[decl]} --- a checked projection over an explicit
proof-carrying record (conditional on the imports and recognition
sources), not an independent Lean derivation. \\
\addlinespace
Theorem, narrowed representation &
Verified in Lean as \texttt{[decl]} over a narrowed representation
(the local syntactic refinement-stable predicate); the narrowing is
recorded below. \\
\addlinespace
Definition, typed structure carrier (import) &
Encoded in Lean as a typed structure carrier \texttt{[decl]} bundling
the imported theorem(s) as hypothesis fields; not a Lean axiom. \\
\addlinespace
Definition, recognition-source carrier &
Supplied as a recognition source, encoded in Lean as a typed structure
carrier \texttt{[decl]} over a narrowed representation; its arithmetic
content is a hypothesis field, not a Lean axiom. \\
\addlinespace
Definition, faithful (typed mirror) &
No footnote; the declaration appears only in the per-label coverage
table below. \\
\addlinespace
No-go / remark / non-claim (not mechanized) &
No footnote and no Lean identifier in the body; recorded as a
paper-level statement only. \\
\bottomrule
\end{tabularx}
\end{center}
\subsection*{Standing narrowings}

\paragraph{Projection-packaged conditional schemas.}
Five results in this paper are mechanized as checked projections over
explicit proof-carrying records rather than as independent Lean
derivations: the \(\chiCTp\) Cassels--Tate comparison, the
overconvergent \(\eta\)-formula, the rank-\(\leq 1\) cascade theorem,
the composite signature, and the headline conditional landing.  This
means that the formalization checks the projection or implication
chain from the supplied records to the stated conclusion.  For example,
in the landing it checks that the recognition sources and headline imports
are threaded through the shell-readout theorems so that the residual
vanishes and the Strong-BSD scalar identity follows.

The corresponding bound is equally important.  These projections do
not independently derive the underlying arithmetic-geometry inputs:
the Cassels--Tate comparison inputs, the overconvergent \(L\)-value
inputs, and the rank-\(\leq 1\) Iwasawa/Heegner inputs remain carried
hypotheses.  This is why the body footnotes for these five results say
that they are tracked as conditional schemas and are not independent
Lean derivations.  In particular, the conditional landing is a theorem
with the recognition sources and headline imports as explicit hypothesis
parameters; it is not an unconditional Lean proof of BSD.

\paragraph{Typed structure carriers.}
The five imported theorem stacks and the three recognition sources are
encoded as typed structures whose arithmetic payload is a hypothesis
field.  Three of the imported stacks, \(\chiCTp\), \(\AofE\), and the
Beilinson comparison, are the headline imports used directly in the landing
theorem; the overconvergent \(\eta\) and \(T_{\mathrm{CASCADE}}\) stacks are
supporting conditional imports recorded by their own sections.  For the
imports, this means that the external theorem stacks are carried as
structured records rather than re-derived.  For the recognition sources, it
means that the named arithmetic content is available as a field of the
carrier and is then consumed by the closure architecture.  In both cases
the content is carried, not derived, and it never appears as a Lean axiom.

The recognition-source carriers also use a narrowed structural
representation.  The carrier records the relevant arithmetic assertion
in typed form---for example the \(p\)-adic descent congruence, the
signed-Selmer persistence comparison, or the higher-rank fixity
readout---without formalizing the full underlying \(p\)-adic, Selmer,
or regulator theory inside Lean.  The dimension and role of each
carrier are faithful to the intended arithmetic content, but the
background arithmetic theory is not redeveloped in the formalization.

Neither of these narrowings enlarges the Lean trust base.  They delimit
what the finite mechanization claims: the imported arithmetic and
recognition content are explicit carried hypotheses, while the checked
Lean layer verifies the structural use of those hypotheses.
\begin{table}[tbp]
\centering
\small
\caption{Standing representation notes for the closure formalization.}
\label{tab:representation-notes}
\begin{tabularx}{\textwidth}{@{}>{\raggedright\arraybackslash}p{0.25\textwidth}
>{\raggedright\arraybackslash}p{0.30\textwidth}
>{\raggedright\arraybackslash}X@{}}
\toprule
Representation issue & Where it appears & Disclosure effect \\
\midrule
Import structures & \(\chiCTp\), \(\AofE\), Beilinson, \(\eta\), and
cascade imports & Arithmetic theorem stacks are hypothesis fields in typed
records, not Lean axioms or internal derivations. \\
\addlinespace
Recognition carriers & \(\GammaPD\), \(\GammaSP\), \(\GammaGZ\) & The
arithmetic payload is supplied closure content and is consumed by the shell
as a hypothesis. \\
\addlinespace
Projection-packaged landing & Supporting conditionals and
\Cref{thm:closure:strong-bsd-conditional} & The formalization checks the
conditional projection/implication chain over supplied records, not an
independent arithmetic proof. \\
\addlinespace
Local \(\RefStable\) predicate & AOR-instance theorem & The Lean theorem
checks the paper-local syntactic predicate, not the full AOR meta-theory. \\
\addlinespace
Cross-paper \(\kapr\) import & \(\Cref{sec:aor-instance}\) & The
normalization is reused from the apparatus paper and cited there, not
re-mechanized in the closure axis. \\
\bottomrule
\end{tabularx}
\end{table}

\subsection*{Non-eliminative AOR instance}

The AOR-instance membership theorem is mechanized over a narrowed
representation.  The development defines, for this paper, a local
syntactic refinement-stable predicate for the BSD trace shell; it does
not mechanize the full Audited Operational Realisability meta-theory of
\cite{TsiokosAOR2026}.  The checked statement is therefore the following
bounded claim: with respect to this local predicate, the declared audit
records of the BSD trace shell account for the structural defects that
persist under refinement.  It is not a claim that Lean has formalized the
general AOR meta-theory, and it is not a claim of membership in that
full meta-theoretic sense.  This is why the body footnote for the
AOR-instance theorem carries the narrowed-representation qualifier
rather than a bare ``verified in Lean'' formulation.

The non-eliminative content is also explicit in the formal record.  The
three recognition sources are discharged as bridged records: each is
recorded in the audit as a primary source-type structural defect, with
its forced secondary role, target, transport, and limit defects, and is
then reclassified as carried recognition content.  None of the three is
marked as eliminated.  By contrast, the mechanical components---the
trace-shell construction, the \(\kapr\)-normalization, and the imported
theorem stacks---are discharged either by construction or by an approved
external-record citation chain.

Thus the AOR-instance membership is non-eliminative.  It records that
the open arithmetic content has been named, carried, and bridged inside
the audit presentation, not that the content has been closed.  In
particular, the rank-\(\geq 2\) generalized Gross--Zagier identity and
the additive \(p \leq 3\) Iwasawa main conjecture remain open after the
AOR reading; the membership records their place in the declared
presentation rather than proving them.  This subsection fixes the exact
scope of the formal AOR result: a local, refinement-stable,
non-eliminative audit-closed membership, so that neither the headline
conditional landing nor the AOR reading is mistaken for an unconditional
or eliminative theorem.
\subsection*{Per-label coverage}

The table lists every closure statement of record in document order,
together with its environment kind, mechanization fidelity, and body
location.  Rows marked faithful, projection, or narrowed have a Lean
declaration, mapped in the next subsection; rows marked not mechanized
are paper-level statements and carry no Lean declaration.

\begin{small}
\begin{longtable}{@{}>{\ttfamily}p{0.45\textwidth} l l l@{}}
\caption{Lean coverage of the labelled statements in the inventory, in document order (30 rows).  Labels omit the \texttt{closure:} infix.}\label{tab:formalization-coverage}\\
\toprule \normalfont Label & Kind & Lean coverage & Section \\ \midrule \endfirsthead
\multicolumn{4}{@{}l}{\footnotesize\itshape continued from previous page}\\ \toprule \normalfont Label & Kind & Lean coverage & Section \\ \midrule \endhead
\midrule \multicolumn{4}{r@{}}{\footnotesize\itshape continued on next page}\\ \endfoot
\bottomrule \endlastfoot
def:chi-ct-p-import & \normalfont definition & \normalfont Lean definition & \normalfont \Cref{sec:imports} \\
def:a-e-import & \normalfont definition & \normalfont Lean definition & \normalfont \Cref{sec:imports} \\
def:beilinson-import & \normalfont definition & \normalfont Lean definition & \normalfont \Cref{sec:imports} \\
nonclaim:imports-not-derived & \normalfont nonclaim & \normalfont not formalized & \normalfont \Cref{sec:imports} \\
def:gamma-padic-descent & \normalfont definition & \normalfont Lean definition & \normalfont \Cref{sec:recognition-sources} \\
def:gamma-sha-persistence & \normalfont definition & \normalfont Lean definition & \normalfont \Cref{sec:recognition-sources} \\
def:gamma-higher-gz-fixity & \normalfont definition & \normalfont Lean definition & \normalfont \Cref{sec:recognition-sources} \\
nonclaim:recognition-sources-not-derived & \normalfont nonclaim & \normalfont not formalized & \normalfont \Cref{sec:recognition-sources} \\
def:sel-bsd-shell & \normalfont definition & \normalfont Lean definition & \normalfont \Cref{sec:shell-readout} \\
def:pi-bsd & \normalfont definition & \normalfont Lean definition & \normalfont \Cref{sec:shell-readout} \\
thm:pi-bsd-iff-strong-bsd & \normalfont theorem & \normalfont Lean proves & \normalfont \Cref{sec:shell-readout} \\
thm:scalar-closure-translation & \normalfont theorem & \normalfont Lean proves & \normalfont \Cref{sec:shell-readout} \\
thm:formed-closure-bridge & \normalfont theorem & \normalfont Lean proves & \normalfont \Cref{sec:shell-readout} \\
thm:composite-signature & \normalfont theorem & \normalfont conditional schema & \normalfont \Cref{sec:shell-readout} \\
rmk:sensitivity-not-necessity & \normalfont remark & \normalfont not formalized & \normalfont \Cref{sec:shell-readout} \\
thm:chi-ct-p-comparison & \normalfont theorem & \normalfont conditional schema & \normalfont \Cref{sec:chi-ct-p} \\
def:eta-published-imports & \normalfont definition & \normalfont Lean definition & \normalfont \Cref{sec:eta-formula} \\
thm:oc-eta-formula & \normalfont theorem & \normalfont conditional schema & \normalfont \Cref{sec:eta-formula} \\
def:t-cascade-import & \normalfont definition & \normalfont Lean definition & \normalfont \Cref{sec:t-cascade} \\
thm:t-cascade-rank-le-1 & \normalfont theorem & \normalfont conditional schema & \normalfont \Cref{sec:t-cascade} \\
rmk:t-e12-gap & \normalfont remark & \normalfont not formalized & \normalfont \Cref{sec:obstructions} \\
rmk:t-bad-gap & \normalfont remark & \normalfont not formalized & \normalfont \Cref{sec:obstructions} \\
rmk:mode-b-residuals & \normalfont remark & \normalfont not formalized & \normalfont \Cref{sec:obstructions} \\
thm:strong-bsd-conditional & \normalfont theorem & \normalfont conditional schema & \normalfont \Cref{sec:landing} \\
nonclaim:landing-conditional & \normalfont nonclaim & \normalfont not formalized & \normalfont \Cref{sec:landing} \\
def:aor-sel-instance & \normalfont definition & \normalfont Lean definition & \normalfont \Cref{sec:aor-instance} \\
thm:aor-mechanical-records & \normalfont theorem & \normalfont model version & \normalfont \Cref{sec:aor-instance} \\
thm:aor-recognition-discharge & \normalfont theorem & \normalfont conditional schema & \normalfont \Cref{sec:aor-instance} \\
thm:aor-instance & \normalfont theorem & \normalfont model version & \normalfont \Cref{sec:aor-instance} \\
rmk:aor-partial-status & \normalfont remark & \normalfont not formalized & \normalfont \Cref{sec:aor-instance} \\
\end{longtable}
\end{small}

\subsection*{Lean declaration map}

The 21 mechanized closure statements each correspond to a Lean
declaration, listed here for citation and audit.  Labels elide the
\texttt{closure:} prefix, and declarations elide the
\texttt{SixBirdsBSD.Closure.} namespace; the seven not-mechanized
statements have no declaration and are omitted from this map, with full
coverage recorded in \Cref{tab:formalization-coverage}.  The
\(\kapr\)-normalization reused in the AOR mechanical-records theorem is
imported from \citet{TsiokosBSDApparatus2026};
its own declaration belongs to that paper rather than to the closure
declaration map.

\begin{small}
\begin{longtable}{@{}>{\ttfamily}p{0.44\textwidth} >{\ttfamily}p{0.50\textwidth}@{}}
\caption{Lean declaration map for the 21 mechanized closure statements. Labels elide the \texttt{closure:} prefix; declarations elide the \texttt{SixBirdsBSD.Closure.} namespace.}\label{tab:appd-declmap-closure}\\
\toprule \normalfont Paper label & \normalfont Lean declaration \\ \midrule \endfirsthead
\multicolumn{2}{@{}l}{\footnotesize\itshape continued from previous page}\\ \toprule \normalfont Paper label & \normalfont Lean declaration \\ \midrule \endhead
\midrule \multicolumn{2}{r@{}}{\footnotesize\itshape continued on next page}\\ \endfoot
\bottomrule \endlastfoot
def:\allowbreak chi-\allowbreak ct-\allowbreak p-\allowbreak import & Imports.\allowbreak chi\allowbreak CTp\allowbreak Import \\
def:\allowbreak a-\allowbreak e-\allowbreak import & Imports.\allowbreak a\allowbreak EImport \\
def:\allowbreak beilinson-\allowbreak import & Imports.\allowbreak beilinson\allowbreak Import \\
def:\allowbreak gamma-\allowbreak padic-\allowbreak descent & Recognition\allowbreak Sources.\allowbreak gamma\allowbreak Padic\allowbreak Descent \\
def:\allowbreak gamma-\allowbreak sha-\allowbreak persistence & Recognition\allowbreak Sources.\allowbreak gamma\allowbreak Sha\allowbreak Persistence \\
def:\allowbreak gamma-\allowbreak higher-\allowbreak gz-\allowbreak fixity & Recognition\allowbreak Sources.\allowbreak gamma\allowbreak Higher\allowbreak GZFixity \\
def:\allowbreak sel-\allowbreak bsd-\allowbreak shell & Sel\allowbreak Shell.\allowbreak sel\allowbreak BSDShell \\
def:\allowbreak pi-\allowbreak bsd & Sel\allowbreak Shell.\allowbreak pi\allowbreak BSD \\
thm:\allowbreak pi-\allowbreak bsd-\allowbreak iff-\allowbreak strong-\allowbreak bsd & Sel\allowbreak Shell.\allowbreak pi\allowbreak BSDIff\allowbreak Strong\allowbreak BSD \\
thm:\allowbreak master-\allowbreak theorem-\allowbreak applicability & Sel\allowbreak Shell.\allowbreak master\allowbreak Theorem\allowbreak Applicability \\
thm:\allowbreak composite-\allowbreak signature & Sel\allowbreak Shell.\allowbreak composite\allowbreak Signature \\
thm:\allowbreak chi-\allowbreak ct-\allowbreak p-\allowbreak comparison & Chi\allowbreak CTp.\allowbreak chi\allowbreak CTp\allowbreak Comparison \\
def:\allowbreak eta-\allowbreak published-\allowbreak imports & Eta\allowbreak Formula.\allowbreak eta\allowbreak Published\allowbreak Imports \\
thm:\allowbreak oc-\allowbreak eta-\allowbreak formula & Eta\allowbreak Formula.\allowbreak oc\allowbreak Eta\allowbreak Formula \\
def:\allowbreak t-\allowbreak cascade-\allowbreak import & TCascade.\allowbreak t\allowbreak Cascade\allowbreak Import \\
thm:\allowbreak t-\allowbreak cascade-\allowbreak rank-\allowbreak le-\allowbreak 1 & TCascade.\allowbreak t\allowbreak Cascade\allowbreak Rank\allowbreak Le\allowbreak One \\
thm:\allowbreak strong-\allowbreak bsd-\allowbreak conditional & Landing.\allowbreak strong\allowbreak BSDConditional \\
def:\allowbreak aor-\allowbreak sel-\allowbreak instance & AORInstance.\allowbreak aor\allowbreak Sel\allowbreak Instance \\
thm:\allowbreak aor-\allowbreak mechanical-\allowbreak records & AORInstance.\allowbreak aor\allowbreak Mechanical\allowbreak Records \\
thm:\allowbreak aor-\allowbreak recognition-\allowbreak discharge & AORInstance.\allowbreak aor\allowbreak Recognition\allowbreak Discharge \\
thm:\allowbreak aor-\allowbreak instance & AORInstance.\allowbreak aor\allowbreak Instance \\
\end{longtable}
\end{small}
